\documentclass[french]{cls/llncs}

\usepackage[english, french]{babel}
\usepackage[utf8]{inputenc}

\usepackage{graphicx}

\usepackage{listings}
\usepackage{subfig}
\usepackage{hyperref}

\title{Rapport projet GNN}

\institute{Université Claude Bernard Lyon 1, France}

\begin{document}
    \maketitle

    \section{Introduction} \label{intro}
    \subsection{Mise en contexte et Problématique}
    La détection d'anomalies classique s'effectue 
    dans des espaces vectoriels euclidiens où les 
    individus sont supposés indépendants. Or, 
    les données réelles (réseaux sociaux, transactions 
    financières, données épidémiologiques) sont 
    fondamentalement interconnectées et modélisées 
    sous forme de graphes.

    Les approches traditionnelles non-profondes 
    (méthodes statistiques, factorisation matricielle) 
    souffrent de limites majeures : elles 
    reposent sur une ingénierie de variables 
    manuelle, passent mal à l'échelle sur de 
    grands réseaux et ne parviennent pas à 
    capturer la non-linéarité des relations et 
    des attributs. L'essor de l'apprentissage 
    profond sur graphes (Graph Representation 
    Learning) vise précisément résoudre ce 
    probleme.

\end{document}